\chapter{Einleitung}

\section{Motivation und Ausgangslage}

Moderne Mikrokern- und Betriebssystemarchitekturen setzen zur Reduzierung der Angriffsfläche konsequent auf das Prinzip der minimalen Rechtevergabe (\textit{Principle of Least Privilege}). Ein zentraler Baustein hierfür sind Capabilities: kryptografisch oder architektonisch unschmiedbare Token, die eine Ressource referenzieren und die darauf zulässigen Operationen fest an die Identität des Tokens binden.

Die vorliegende Arbeit baut direkt auf der Masterarbeit von Julian Schnock auf, deren Entwicklungsstand aus September 2025 datiert. In dieser Vorarbeit wurde ein grundlegendes Capability-System für das Betriebssystem implementiert, welches auf typisierten Capabilities sowie task-lokalen Capability Spaces (CSpaces) basiert. Obwohl diese Basisarchitektur die grundlegende Machbarkeit nachwies, wies der Stand von September 2025 signifikante Schwachstellen hinsichtlich Nebenläufigkeit und Fehlerbehandlung auf. Zudem existierte eine erhebliche Versionsdiskrepanz zum weiterentwickelten Hauptzweig des Kernels.

Das Ziel dieser Arbeit ist die Beseitigung bestehender Synchronisations- und Stabilitätsprobleme sowie die Zusammenführung (Merge) des Capability-Zweigs mit dem aktuellen Kernel-Stand als Fundament für die Integration des Messaging-Systems und von Shared Memory.